\icode{N18J-MG-SI-L3-261021-010}
\nceeID[6]{2018北京}{17}
若想检验“使月球绕地球运动的力”与“使苹果落地的力”遵循同样的规律，在已知月地距离约为地球半径60倍的情况下，需要验证\paren[]
\begin{choices}
	\item 地球吸引月球的力约为地球吸引苹果的力的\(\frac{1}{60^2}\)
	\item 月球公转的加速度约为苹果落向地面加速度的\(\frac{1}{60^2}\)
	\item 自由落体在月球表面的加速度约为地球表面的\(\frac{1}{6}\)
	\item 苹果在月球表面受到的引力约为在地球表面的\(\frac{1}{60}\)
\end{choices}
\begin{answer}
	B\qquad 提示：万有引力提供月球做圆周运动的向心力，在月球表面的物体受到的万有引力等于重力，据此求出月球表面的重力加速度，从而即可求解
\end{answer}
